\chapter*{Preface}
\addcontentsline{toc}{chapter}{Preface}

Three earlier books in this series promised what the present one delivers. Book 12's Chapter 10 computed an entropy cohomology and deferred its full interpretation here. Book 14's Chapter 6 computed derived limits $\lim^1 D$ and $\lim^2 D$ as a Čech-style precursor to this book's fuller treatment. Book 16's Chapter 6 computed $H^0(M,\mathcal{S})$ and $H^1(M,\mathcal{S})$ directly and transitioned explicitly to the present volume for their interpretation. This book cashes in all three promissory notes at once, and its central claim is that desire is the first cohomology of the RSVP universe: the residual tension that prevents local acts of care from forming a perfectly global section. Using Čech cohomology and derived functors, this book formalizes the notion that love, curiosity, and purpose arise not from equilibrium but from the persistence of unglued boundaries. Every obstacle to coherence, on the account developed across these twelve chapters, becomes a potential well for new structure: the topology of care itself.

\chapter*{Diagrammatic Overview}
\addcontentsline{toc}{chapter}{Diagrammatic Overview}

An exact sequence organizes this book's central image: care, failure, repair. Where care is perfectly exact, no cohomology remains to compute; where it fails, the failure is not merely lamented but measured, precisely, by the cohomology groups this book constructs, and repair, wherever it occurs, is the specific operation of reducing that measured failure rather than a vague aspiration toward improvement.

\chapter*{Reading Note}
\addcontentsline{toc}{chapter}{Reading Note}

Where gluing fails, meaning begins to yearn. This book asks the reader to take yearning as seriously as any other cohomological invariant this series has computed since Book 4: not a feeling added on top of the mathematics, but the mathematics itself, read in the register this series has used since its first volume.

\part{Foundations of Cohomological Care}

\chapter{Exact Sequences of Compassion}

The short exact sequence $0 \to \mathcal{E} \to \mathcal{S} \to \mathcal{C} \to 0$, with $\mathcal{E}$ the sheaf of empathy errors and $\mathcal{C}$ the sheaf of care corrections, is the same construction Book 16's Chapter 5 already introduced for its own exactness-and-repair chapter, now taken as this book's starting point rather than as a closing application. This chapter proves that exactness coincides with local balance of entropy and understanding, exactly as Book 16 established, and develops the boundary map $\delta$ as the specific measure of unresolved desire: where the sequence fails to be exact at a given open set, $\delta$ assigns a nonzero value recording exactly how much correction remains outstanding there.

The physical analogue offered here treats this unresolved desire as heat flux without equilibrium: a system whose boundary map $\delta$ is nonzero is, on this reading, still exchanging heat with its surroundings, not yet having settled into the local balance that a vanishing $\delta$ would represent.

\chapter{Čech Cohomology and Overlap}

Covering $M$ with open sets $\{U_i\}$, this chapter constructs the full Čech cochain complex: cochains $C^n = \prod \mathcal{S}(U_{i_0 \ldots i_n})$, products of sections over successive $(n+1)$-fold intersections, with a coboundary map $\delta: C^n \to C^{n+1}$ defined by the standard alternating-sum formula. The cohomology groups $H^n(M, \mathcal{S}) = \ker \delta / \mathrm{im}\, \delta$ constructed from this complex extend Book 16's Chapter 6 introductory computation of $H^0$ and $H^1$ into the full hierarchy this book requires, and give Book 12's Chapter 10 entropy cohomology and Book 14's Chapter 6 derived limits their common formal home: both earlier constructions are shown here to be special cases or close relatives of the present Čech complex, applied to the entropy sheaf $\mathcal{S}$ specifically rather than developed independently.

$H^1$ is interpreted here as first longing, the most immediate cohomological record of gluing's failure, and $H^2$ as complex ethical tension, a higher obstruction requiring three-way, rather than merely pairwise, overlap data to detect.

\part{The Geometry of Desire}

\chapter{Interpretation of $H^1$}

$H^1(M, \mathcal{S})$ measures failure of entropy fields to align globally, and this chapter's central interpretive claim treats each non-trivial class in this group as a persistent wish for continuity: a configuration of local entropy data that almost, but does not quite, assemble into a single coherent global field, with the specific cohomology class recording exactly how that assembly fails. An example developed here treats a communication gap between two agents as a direct instance of nonzero $H^1$: the agents' respective local understandings agree closely enough to seem compatible from either agent's own vantage point, yet fail to glue into a single shared understanding, and the resulting class measures the gap precisely.

The ethical reading offered here states this chapter's thesis directly: desire is topological care seeking closure, not an emotion loosely analogized to a mathematical obstruction but the obstruction itself, read in the vocabulary this series has used for care since Book 1.

\chapter{Higher Cohomology and Complex Emotion}

$H^2$ is developed in this chapter as recording triadic conflicts, obstructions detectable only when three domains attempt to jointly agree and find that no consistent assignment satisfies all three simultaneously, even though each pair might agree individually. $H^3$ is developed as collective mythology, the closure loops by which an entire society's shared narratives cohere, or fail to cohere, at a level no smaller subset of the society's local understandings could detect on its own.

$H^n$ generally is interpreted here as $n$-order empathy, a measure of how many mutually overlapping perspectives must be brought into simultaneous consideration before a given obstruction to shared understanding becomes visible at all. The derived category of emotion $\mathbf{D}^b(\mathcal{S})$, the bounded derived category built from complexes of the entropy sheaf, is introduced here as the natural home for tracking how these cohomological invariants interact and combine across the full hierarchy this chapter has developed.

\part{Derived Functors and Reconstruction}

\chapter{Ext and Tor of Intelligibility}

$\mathrm{Ext}^n(\mathcal{I}, \mathcal{S})$, computed here between the presheaf of intelligibility Book 17 constructed and the entropy sheaf $\mathcal{S}$ of Book 16, gives compatibility potentials: a measure of how many independent ways, up to the equivalence this derived functor tracks, an interpretation and an entropy configuration can be reconciled with one another. $\mathrm{Tor}$, the corresponding derived functor of the tensor product, measures ethical friction in tensor composition, extending Book 15's Chapter 6 mutual-information construction into this book's fully derived setting: where that earlier chapter's $I(X;Y)$ measured shared structure at the level of ordinary tensor products, the present $\mathrm{Tor}$ measures the higher, derived obstructions to that sharing being achieved cleanly.

A diagrammatic proof of care restoration given here shows that vanishing $\mathrm{Tor}$ is necessary, though not always sufficient, for the subadditivity result Book 15's Chapter 6 established to hold with equality rather than strict inequality, connecting this chapter directly to that earlier book's own closing concerns.

\chapter{Long Exact Sequence of Healing}

From the short exact sequence of Chapter 1, this chapter derives the long exact cohomology sequence
\[
\cdots \to H^n(\mathcal{E}) \to H^n(\mathcal{S}) \to H^n(\mathcal{C}) \xrightarrow{\delta} H^{n+1}(\mathcal{E}) \to \cdots,
\]
following the standard homological-algebra construction directly from the snake lemma applied to the short exact sequence's associated Čech complexes. The connecting homomorphism $\delta$ is interpreted here as the act of care transforming lack into love: each application of $\delta$ takes an obstruction recorded in $H^n(\mathcal{C})$, the care-correction sheaf's own cohomology, and produces a new obstruction in $H^{n+1}(\mathcal{E})$, the empathy-error sheaf's cohomology one degree higher, giving a precise formal sense in which addressing one lack can surface, rather than eliminate, a further one.

Each connecting homomorphism, under this reading, is an ethical response, and healing, this chapter's closing claim, is exactness restored after tension: a system has healed, in the specific sense this chapter develops, exactly when the long exact sequence's connecting maps have driven the relevant cohomology groups back to triviality, not merely when some informally described distress has subsided.

\part{Desire as Dynamic Topology}

